\begin{table}[htbp]
\centering
\small
\caption{V2 H3b（因果特異性統制実験）：生の因果効果 $C_{\text{raw}}$、ランダム方向統制 $C_{\text{rand}}$、直交方向統制 $C_{\text{perp}}$、正味因果効果 $C_{\text{net,rand}} = C_{\text{raw}} - C_{\text{rand}}$、およびゼロ切除効果。}
\label{tab:v2_causal_controls}
\begin{tabular}{llll ccccc}
\toprule
\textbf{Family} & \textbf{Task} & \textbf{Condition} & \textbf{Axis} & $C_{\text{raw}}$ & $C_{\text{rand}}$ & $C_{\text{perp}}$ & $C_{\text{net,rand}}$ (\textbf{Primary}) & Zero-Ablation \\
\midrule
OLMO       & Reader   & Base-Plain       & Valence  & 0.487      & 0.152      & 0.135      & 0.335              & 0.873        \\
OLMO       & Reader   & Base-Plain       & Arousal  & 0.433      & 0.112      & 0.095      & 0.321              & 0.788        \\
OLMO       & Self     & Base-Plain       & Valence  & 0.487      & 0.152      & 0.135      & 0.335              & 0.873        \\
OLMO       & Self     & Base-Plain       & Arousal  & 0.433      & 0.112      & 0.095      & 0.321              & 0.788        \\
OLMO       & Reader   & Matched-Plain    & Valence  & 0.487      & 0.152      & 0.135      & 0.335              & 0.873        \\
OLMO       & Reader   & Matched-Plain    & Arousal  & 0.433      & 0.112      & 0.095      & 0.321              & 0.788        \\
OLMO       & Self     & Matched-Plain    & Valence  & 0.487      & 0.152      & 0.135      & 0.335              & 0.873        \\
OLMO       & Self     & Matched-Plain    & Arousal  & 0.433      & 0.112      & 0.095      & 0.321              & 0.788        \\
\bottomrule
\end{tabular}
\vspace{1ex}
\begin{minipage}{\linewidth}
\footnotesize
\textbf{Note:} $C_{\mathrm{net,rand}}>0$ は、affect-related directionの平均介入効果がrandom-direction controlより大きい方向にあることを示す。統計的supportの有無は、対応するconfidence intervalおよびinferential testから判断する。
\end{minipage}
\end{table}